\chapter{Render Pipeline Guide}\label{record-014}

\noindent{\footnotesize\color{BookMuted}\path{docs/guides/render-pipeline.md}\par}\smallskip

The render pipeline in \texttt{@\allowbreak{}f\allowbreak{}l\allowbreak{}y\allowbreak{}i\allowbreak{}n\allowbreak{}g\allowbreak{}r\allowbreak{}o\allowbreak{}b\allowbreak{}o\allowbreak{}t\allowbreak{}s\allowbreak{}/\allowbreak{}b\allowbreak{}i\allowbreak{}j\allowbreak{}o\allowbreak{}u\allowbreak{}-\allowbreak{}t\allowbreak{}u\allowbreak{}i} is a deterministic, stage-ordered execution chain.

It runs as:

\begin{enumerate}
\def\labelenumi{\arabic{enumi}.}
\tightlist
\item
  \texttt{L\allowbreak{}a\allowbreak{}y\allowbreak{}o\allowbreak{}u\allowbreak{}t}
\item
  \texttt{P\allowbreak{}a\allowbreak{}i\allowbreak{}n\allowbreak{}t}
\item
  \texttt{P\allowbreak{}o\allowbreak{}s\allowbreak{}t\allowbreak{}P\allowbreak{}r\allowbreak{}o\allowbreak{}c\allowbreak{}e\allowbreak{}s\allowbreak{}s}
\item
  \texttt{D\allowbreak{}i\allowbreak{}f\allowbreak{}f}
\item
  \texttt{O\allowbreak{}u\allowbreak{}t\allowbreak{}p\allowbreak{}u\allowbreak{}t}
\end{enumerate}

\section{Default Runtime Wiring}\label{record-014-default-runtime-wiring}

\texttt{r\allowbreak{}u\allowbreak{}n\allowbreak{}(\allowbreak{})} installs middleware in this fixed stage order:

\begin{enumerate}
\def\labelenumi{\arabic{enumi}.}
\tightlist
\item
  \texttt{L\allowbreak{}a\allowbreak{}y\allowbreak{}o\allowbreak{}u\allowbreak{}t}: call \texttt{a\allowbreak{}p\allowbreak{}p\allowbreak{}.\allowbreak{}v\allowbreak{}i\allowbreak{}e\allowbreak{}w\allowbreak{}(\allowbreak{}m\allowbreak{}o\allowbreak{}d\allowbreak{}e\allowbreak{}l\allowbreak{})} and produce the layout output.
\item
  \texttt{L\allowbreak{}a\allowbreak{}y\allowbreak{}o\allowbreak{}u\allowbreak{}t}: apply motion reconciliation.
\item
  \texttt{L\allowbreak{}a\allowbreak{}y\allowbreak{}o\allowbreak{}u\allowbreak{}t}: apply BCSS middleware when \texttt{o\allowbreak{}p\allowbreak{}t\allowbreak{}i\allowbreak{}o\allowbreak{}n\allowbreak{}s\allowbreak{}.\allowbreak{}c\allowbreak{}s\allowbreak{}s} is present.
\item
  \texttt{P\allowbreak{}a\allowbreak{}i\allowbreak{}n\allowbreak{}t}: paint the resolved layout tree into \texttt{t\allowbreak{}a\allowbreak{}r\allowbreak{}g\allowbreak{}e\allowbreak{}t\allowbreak{}S\allowbreak{}u\allowbreak{}r\allowbreak{}f\allowbreak{}a\allowbreak{}c\allowbreak{}e}.
\item
  \texttt{D\allowbreak{}i\allowbreak{}f\allowbreak{}f}: compute terminal diff and fill \texttt{o\allowbreak{}u\allowbreak{}t\allowbreak{}B\allowbreak{}u\allowbreak{}f}.
\item
  \texttt{O\allowbreak{}u\allowbreak{}t\allowbreak{}p\allowbreak{}u\allowbreak{}t}: emit diff and swap framebuffers.
\item
  \texttt{c\allowbreak{}o\allowbreak{}n\allowbreak{}f\allowbreak{}i\allowbreak{}g\allowbreak{}u\allowbreak{}r\allowbreak{}e\allowbreak{}P\allowbreak{}i\allowbreak{}p\allowbreak{}e\allowbreak{}l\allowbreak{}i\allowbreak{}n\allowbreak{}e\allowbreak{}(\allowbreak{}p\allowbreak{}i\allowbreak{}p\allowbreak{}e\allowbreak{}l\allowbreak{}i\allowbreak{}n\allowbreak{}e\allowbreak{})} if provided by the caller.
\end{enumerate}

Custom middleware runs after the built-ins in the same stage and cannot change stage ordering.

\section{Default and Custom Stage Contract}\label{record-014-default-and-custom-stage-contract}

Use \texttt{c\allowbreak{}o\allowbreak{}n\allowbreak{}f\allowbreak{}i\allowbreak{}g\allowbreak{}u\allowbreak{}r\allowbreak{}e\allowbreak{}P\allowbreak{}i\allowbreak{}p\allowbreak{}e\allowbreak{}l\allowbreak{}i\allowbreak{}n\allowbreak{}e\allowbreak{}(\allowbreak{})} when your middleware belongs to one of these five stages:

\begin{itemize}
\tightlist
\item
  \texttt{L\allowbreak{}a\allowbreak{}y\allowbreak{}o\allowbreak{}u\allowbreak{}t}: inspect or rewrite \texttt{L\allowbreak{}a\allowbreak{}y\allowbreak{}o\allowbreak{}u\allowbreak{}t\allowbreak{}N\allowbreak{}o\allowbreak{}d\allowbreak{}e} trees.
\item
  \texttt{P\allowbreak{}a\allowbreak{}i\allowbreak{}n\allowbreak{}t}: mutate \texttt{t\allowbreak{}a\allowbreak{}r\allowbreak{}g\allowbreak{}e\allowbreak{}t\allowbreak{}S\allowbreak{}u\allowbreak{}r\allowbreak{}f\allowbreak{}a\allowbreak{}c\allowbreak{}e} by painting extra marks directly.
\item
  \texttt{P\allowbreak{}o\allowbreak{}s\allowbreak{}t\allowbreak{}P\allowbreak{}r\allowbreak{}o\allowbreak{}c\allowbreak{}e\allowbreak{}s\allowbreak{}s}: apply surface-wide shaders or instrumentation before diff.
\item
  \texttt{D\allowbreak{}i\allowbreak{}f\allowbreak{}f}: adjust output encoding just before emit (advanced usage).
\item
  \texttt{O\allowbreak{}u\allowbreak{}t\allowbreak{}p\allowbreak{}u\allowbreak{}t}: run end-of-frame side-effects after terminal write and swap.
\end{itemize}

\section{Inter-Stage Data Contracts}\label{record-014-inter-stage-data-contracts}

A \texttt{R\allowbreak{}e\allowbreak{}n\allowbreak{}d\allowbreak{}e\allowbreak{}r\allowbreak{}S\allowbreak{}t\allowbreak{}a\allowbreak{}t\allowbreak{}e} instance is shared across all middleware for a single frame. The key shared state channels are:

\begin{itemize}
\tightlist
\item
  \texttt{m\allowbreak{}o\allowbreak{}d\allowbreak{}e\allowbreak{}l}: current model value for this frame.
\item
  \texttt{l\allowbreak{}a\allowbreak{}y\allowbreak{}o\allowbreak{}u\allowbreak{}t\allowbreak{}M\allowbreak{}a\allowbreak{}p}: geometry metadata for resolved nodes.
\item
  \texttt{d\allowbreak{}a\allowbreak{}t\allowbreak{}a}: unstructured scratchbag intended for middleware communication.
\end{itemize}

\subsection{\texorpdfstring{\texttt{l\allowbreak{}a\allowbreak{}y\allowbreak{}o\allowbreak{}u\allowbreak{}t\allowbreak{}R\allowbreak{}o\allowbreak{}o\allowbreak{}t} via \texttt{R\allowbreak{}e\allowbreak{}n\allowbreak{}d\allowbreak{}e\allowbreak{}r\allowbreak{}S\allowbreak{}t\allowbreak{}a\allowbreak{}t\allowbreak{}e\allowbreak{}.\allowbreak{}d\allowbreak{}a\allowbreak{}t\allowbreak{}a}}{ via }}\label{record-014-layoutroot-via-renderstate.data}

In current runtime code, the \texttt{L\allowbreak{}a\allowbreak{}y\allowbreak{}o\allowbreak{}u\allowbreak{}t} stage places the resolved layout tree in \texttt{s\allowbreak{}t\allowbreak{}a\allowbreak{}t\allowbreak{}e\allowbreak{}.\allowbreak{}d\allowbreak{}a\allowbreak{}t\allowbreak{}a\allowbreak{}[\allowbreak{}'\allowbreak{}\_\allowbreak{}\_\allowbreak{}b\allowbreak{}i\allowbreak{}j\allowbreak{}o\allowbreak{}u\allowbreak{}:\allowbreak{}l\allowbreak{}a\allowbreak{}y\allowbreak{}o\allowbreak{}u\allowbreak{}t\allowbreak{}R\allowbreak{}o\allowbreak{}o\allowbreak{}t\allowbreak{}'\allowbreak{}]}.

That object is then consumed by \texttt{m\allowbreak{}o\allowbreak{}t\allowbreak{}i\allowbreak{}o\allowbreak{}n\allowbreak{}M\allowbreak{}i\allowbreak{}d\allowbreak{}d\allowbreak{}l\allowbreak{}e\allowbreak{}w\allowbreak{}a\allowbreak{}r\allowbreak{}e\allowbreak{}(\allowbreak{})} and \texttt{p\allowbreak{}a\allowbreak{}i\allowbreak{}n\allowbreak{}t\allowbreak{}M\allowbreak{}i\allowbreak{}d\allowbreak{}d\allowbreak{}l\allowbreak{}e\allowbreak{}w\allowbreak{}a\allowbreak{}r\allowbreak{}e\allowbreak{}(\allowbreak{})} in later stages.

Recommended access pattern for custom middleware:

\begin{itemize}
\tightlist
\item
  Read and write your own namespaced keys, for example:

  \begin{itemize}
  \tightlist
  \item
    \texttt{s\allowbreak{}t\allowbreak{}a\allowbreak{}t\allowbreak{}e\allowbreak{}.\allowbreak{}d\allowbreak{}a\allowbreak{}t\allowbreak{}a\allowbreak{}[\allowbreak{}'\allowbreak{}m\allowbreak{}y\allowbreak{}-\allowbreak{}a\allowbreak{}p\allowbreak{}p\allowbreak{}:\allowbreak{}f\allowbreak{}o\allowbreak{}c\allowbreak{}u\allowbreak{}s\allowbreak{}S\allowbreak{}t\allowbreak{}a\allowbreak{}t\allowbreak{}e\allowbreak{}'\allowbreak{}]}
  \item
    \texttt{s\allowbreak{}t\allowbreak{}a\allowbreak{}t\allowbreak{}e\allowbreak{}.\allowbreak{}d\allowbreak{}a\allowbreak{}t\allowbreak{}a\allowbreak{}[\allowbreak{}'\allowbreak{}m\allowbreak{}y\allowbreak{}-\allowbreak{}a\allowbreak{}p\allowbreak{}p\allowbreak{}:\allowbreak{}t\allowbreak{}r\allowbreak{}a\allowbreak{}c\allowbreak{}e\allowbreak{}:\allowbreak{}e\allowbreak{}n\allowbreak{}a\allowbreak{}b\allowbreak{}l\allowbreak{}e\allowbreak{}d\allowbreak{}'\allowbreak{}]}
  \end{itemize}
\end{itemize}

This avoids private runtime keys and avoids coupling across feature layers.

The built-in layout keys are intentionally not part of your public app contract.

\subsection{LayoutMap Semantics}\label{record-014-layoutmap-semantics}

\texttt{l\allowbreak{}a\allowbreak{}y\allowbreak{}o\allowbreak{}u\allowbreak{}t\allowbreak{}M\allowbreak{}a\allowbreak{}p} is currently populated with keyed geometry for internal systems such as a11y overlays, layout debugging, and event routing.

If you need layout geometry from custom middleware, guard key lookups and treat missing values as an expected pre-layout state.

\section{Framebuffer Flip Model (Current Surface Truth)}\label{record-014-framebuffer-flip-model-current-surface-truth}

\texttt{D\allowbreak{}i\allowbreak{}f\allowbreak{}f} compares:

\begin{itemize}
\tightlist
\item
  \texttt{s\allowbreak{}t\allowbreak{}a\allowbreak{}t\allowbreak{}e\allowbreak{}.\allowbreak{}c\allowbreak{}u\allowbreak{}r\allowbreak{}r\allowbreak{}e\allowbreak{}n\allowbreak{}t\allowbreak{}S\allowbreak{}u\allowbreak{}r\allowbreak{}f\allowbreak{}a\allowbreak{}c\allowbreak{}e}: the previous visible surface
\item
  \texttt{s\allowbreak{}t\allowbreak{}a\allowbreak{}t\allowbreak{}e\allowbreak{}.\allowbreak{}t\allowbreak{}a\allowbreak{}r\allowbreak{}g\allowbreak{}e\allowbreak{}t\allowbreak{}S\allowbreak{}u\allowbreak{}r\allowbreak{}f\allowbreak{}a\allowbreak{}c\allowbreak{}e}: the next surface under construction
\end{itemize}

After \texttt{O\allowbreak{}u\allowbreak{}t\allowbreak{}p\allowbreak{}u\allowbreak{}t} runs, the runtime flips the internal buffers.

The conceptual model is:

\begin{CodeBlock}
before Output:   currentSurface is frame N, targetSurface is frame N+1
after Output:    currentSurface is frame N+1, targetSurface becomes previous frame N
\end{CodeBlock}

That means middleware appended to \texttt{O\allowbreak{}u\allowbreak{}t\allowbreak{}p\allowbreak{}u\allowbreak{}t} after swap must not assume that \texttt{s\allowbreak{}t\allowbreak{}a\allowbreak{}t\allowbreak{}e\allowbreak{}.\allowbreak{}c\allowbreak{}u\allowbreak{}r\allowbreak{}r\allowbreak{}e\allowbreak{}n\allowbreak{}t\allowbreak{}S\allowbreak{}u\allowbreak{}r\allowbreak{}f\allowbreak{}a\allowbreak{}c\allowbreak{}e} still points at the old frame unless it reads it before invoking default output semantics.

In practical terms:

\begin{itemize}
\tightlist
\item
  Use \texttt{c\allowbreak{}u\allowbreak{}r\allowbreak{}r\allowbreak{}e\allowbreak{}n\allowbreak{}t\allowbreak{}S\allowbreak{}u\allowbreak{}r\allowbreak{}f\allowbreak{}a\allowbreak{}c\allowbreak{}e} for diff inputs and visibility of the previous frame.
\item
  Use \texttt{t\allowbreak{}a\allowbreak{}r\allowbreak{}g\allowbreak{}e\allowbreak{}t\allowbreak{}S\allowbreak{}u\allowbreak{}r\allowbreak{}f\allowbreak{}a\allowbreak{}c\allowbreak{}e} as the paint target for the frame currently being built.
\item
  Do not retain references across frames outside \texttt{L\allowbreak{}a\allowbreak{}y\allowbreak{}o\allowbreak{}u\allowbreak{}t}/\texttt{P\allowbreak{}a\allowbreak{}i\allowbreak{}n\allowbreak{}t} middleware.
\end{itemize}

\section{Why the Pipeline Is Synchronous}\label{record-014-why-the-pipeline-is-synchronous}

The pipeline is synchronous by design to keep frame cost measurable and predictable.

Bijou's TEA runtime calls \texttt{v\allowbreak{}i\allowbreak{}e\allowbreak{}w\allowbreak{}(\allowbreak{}m\allowbreak{}o\allowbreak{}d\allowbreak{}e\allowbreak{}l\allowbreak{})} for a frame. That full-model render is intentional: state is explicit, the frame can be reproduced from one model snapshot, and tests can ask ``what would this model render?'' without replaying hidden mutations.

The tradeoff is that expensive view work is paid on every rendered frame. Keep \texttt{v\allowbreak{}i\allowbreak{}e\allowbreak{}w\allowbreak{}(\allowbreak{})} pure and bounded:

\begin{itemize}
\tightlist
\item
  store durable app facts in the model,
\item
  derive only the visible slice needed for the current frame,
\item
  cache expensive non-visual preparation outside \texttt{v\allowbreak{}i\allowbreak{}e\allowbreak{}w\allowbreak{}(\allowbreak{})} when the model changes,
\item
  use commands for IO, timers, and other async work.
\end{itemize}

If each stage can run with bounded synchronous cost, the frame time budget is:

\begin{CodeBlock}
frameCost = layout + paint + postProcess + diff + output + middlewareOverhead
\end{CodeBlock}

If an asynchronous stage is inserted without careful partitioning, total frame cost becomes:

\begin{CodeBlock}
frameCost = layout + paint + postProcess + await(someAsyncIO) + diff + output
\end{CodeBlock}

That removes a strict per-frame budget and can stall rendering on unrelated IO.

The architectural rule is therefore:

\begin{itemize}
\tightlist
\item
  Keep rendering synchronous and deterministic.
\item
  Move async work into commands.
\item
  Receive async results as messages, then update model and re-render.
\end{itemize}

Runtime diagnostics enforce that posture. If a render middleware returns a Promise, Bijou writes a pipeline diagnostic, keeps the current frame moving, and ignores any late \texttt{n\allowbreak{}e\allowbreak{}x\allowbreak{}t\allowbreak{}(\allowbreak{})} call for stage ordering. A rejected async middleware Promise is also reported through the pipeline error channel. The diagnostic is not a supported async render mode; it is a guardrail for finding accidental \texttt{a\allowbreak{}s\allowbreak{}y\allowbreak{}n\allowbreak{}c} middleware before it turns into frame latency or out-of-order painting.

\section{Worked Example}\label{record-014-worked-example}

The runtime pipeline supports middleware insertion points that preserve the fixed stage ordering.

\begin{CodeBlock}
import { run } from '@flyingrobots/bijou-tui';

run(app, {
  configurePipeline(pipeline) {
    pipeline.use('PostProcess', (state, next) => {
      state.targetSurface.set(0, 0, { char: 'X', empty: false });
      next();
    });
  },
});
\end{CodeBlock}

Because this middleware is registered before app execution, it runs on every frame after all built-ins for \texttt{P\allowbreak{}o\allowbreak{}s\allowbreak{}t\allowbreak{}P\allowbreak{}r\allowbreak{}o\allowbreak{}c\allowbreak{}e\allowbreak{}s\allowbreak{}s} and before default diffing.

\section{Timing Hooks}\label{record-014-timing-hooks}

You can observe per-stage timings from middleware.

\begin{CodeBlock}
import { getRenderStageTimings, run } from '@flyingrobots/bijou-tui';

run(app, {
  configurePipeline(pipeline) {
    pipeline.onStageComplete((stage, durationMs, state) => {
      console.log(stage, durationMs);
      console.log(getRenderStageTimings(state));
    });
  },
});
\end{CodeBlock}

Notes:

\begin{itemize}
\tightlist
\item
  \texttt{o\allowbreak{}n\allowbreak{}S\allowbreak{}t\allowbreak{}a\allowbreak{}g\allowbreak{}e\allowbreak{}C\allowbreak{}o\allowbreak{}m\allowbreak{}p\allowbreak{}l\allowbreak{}e\allowbreak{}t\allowbreak{}e\allowbreak{}(\allowbreak{}.\allowbreak{}.\allowbreak{}.\allowbreak{})} fires once per stage, in fixed stage order.
\item
  \texttt{g\allowbreak{}e\allowbreak{}t\allowbreak{}R\allowbreak{}e\allowbreak{}n\allowbreak{}d\allowbreak{}e\allowbreak{}r\allowbreak{}S\allowbreak{}t\allowbreak{}a\allowbreak{}g\allowbreak{}e\allowbreak{}T\allowbreak{}i\allowbreak{}m\allowbreak{}i\allowbreak{}n\allowbreak{}g\allowbreak{}s\allowbreak{}(\allowbreak{}s\allowbreak{}t\allowbreak{}a\allowbreak{}t\allowbreak{}e\allowbreak{})} returns timings captured for that frame.
\item
  The timings array is replaced on every \texttt{p\allowbreak{}i\allowbreak{}p\allowbreak{}e\allowbreak{}l\allowbreak{}i\allowbreak{}n\allowbreak{}e\allowbreak{}.\allowbreak{}e\allowbreak{}x\allowbreak{}e\allowbreak{}c\allowbreak{}u\allowbreak{}t\allowbreak{}e\allowbreak{}(\allowbreak{}s\allowbreak{}t\allowbreak{}a\allowbreak{}t\allowbreak{}e\allowbreak{})} call.
\item
  Empty stages are represented in timing output.
\end{itemize}

\section{\texorpdfstring{\texttt{R\allowbreak{}e\allowbreak{}n\allowbreak{}d\allowbreak{}e\allowbreak{}r\allowbreak{}S\allowbreak{}t\allowbreak{}a\allowbreak{}t\allowbreak{}e} Reference}{ Reference}}\label{record-014-renderstate-reference}

The public shape is:

\begin{itemize}
\tightlist
\item
  \texttt{m\allowbreak{}o\allowbreak{}d\allowbreak{}e\allowbreak{}l}: read by convention in rendering.
\item
  \texttt{c\allowbreak{}t\allowbreak{}x}: active \texttt{B\allowbreak{}i\allowbreak{}j\allowbreak{}o\allowbreak{}u\allowbreak{}C\allowbreak{}o\allowbreak{}n\allowbreak{}t\allowbreak{}e\allowbreak{}x\allowbreak{}t}.
\item
  \texttt{d\allowbreak{}t}: wall time delta in seconds.
\item
  \texttt{c\allowbreak{}u\allowbreak{}r\allowbreak{}r\allowbreak{}e\allowbreak{}n\allowbreak{}t\allowbreak{}S\allowbreak{}u\allowbreak{}r\allowbreak{}f\allowbreak{}a\allowbreak{}c\allowbreak{}e}: previous frame surface, normally treated as read-only.
\item
  \texttt{t\allowbreak{}a\allowbreak{}r\allowbreak{}g\allowbreak{}e\allowbreak{}t\allowbreak{}S\allowbreak{}u\allowbreak{}r\allowbreak{}f\allowbreak{}a\allowbreak{}c\allowbreak{}e}: current frame paint surface.
\item
  \texttt{o\allowbreak{}u\allowbreak{}t\allowbreak{}B\allowbreak{}u\allowbreak{}f}: optional UTF-8 output buffer for default diff stage.
\item
  \texttt{l\allowbreak{}a\allowbreak{}y\allowbreak{}o\allowbreak{}u\allowbreak{}t\allowbreak{}M\allowbreak{}a\allowbreak{}p}: geometry metadata map.
\item
  \texttt{d\allowbreak{}a\allowbreak{}t\allowbreak{}a}: middleware coordination bag.
\end{itemize}

\texttt{d\allowbreak{}a\allowbreak{}t\allowbreak{}a} is also where layout root and per-frame timing metadata are held.

\section{Stage Failure Semantics}\label{record-014-stage-failure-semantics}

A middleware exception does not kill the app immediately.

\begin{itemize}
\tightlist
\item
  The exception is routed through the same diagnostics path as other runtime faults.
\item
  The frame may render partially or degrade for that tick.
\item
  Rendering continues so the app remains recoverable and operational.
\end{itemize}

If your custom middleware must fail fast, emit a command or message to transition into an intentional error state during \texttt{u\allowbreak{}p\allowbreak{}d\allowbreak{}a\allowbreak{}t\allowbreak{}e\allowbreak{}(\allowbreak{})}.

\section{Read Next}\label{record-014-read-next}

\begin{itemize}
\tightlist
\item
  \href{https://github.com/flyingrobots/bijou/blob/cbd70612140e7a450d764b3c7be1c6af878b99bc/docs/strategy/layout-localization-pipeline.md}{Layout Localization Pipeline}
\item
  \href{https://github.com/flyingrobots/bijou/blob/cbd70612140e7a450d764b3c7be1c6af878b99bc/packages/bijou-tui/ADVANCED_GUIDE.md}{packages/bijou-tui/ADVANCED\_GUIDE.md}
\end{itemize}

